\section{The small negative zeros}

Throughout this section $n\ge250$, $p=n+1$ and $K=\lfloor\sqrt n/3\rfloor$. Set
\[
c(t)=\frac{t-1}{\log t}\quad(t>1),\qquad c(1)=1,\qquad \ell(t)=\log c(t).
\]

\subsection{Coefficient estimates}
For $t\ge1$,
\begin{equation}\label{eq:logmean-curvature}
\ell'(t)\ge\frac1{2t},\qquad -\ell''(t)\ge\frac5{12t^2}.
\end{equation}
For the first inequality, put $u=\log t$ and use
$t\ell'(t)=1/(1-e^{-u})-1/u\ge1/2$.
For the second, multiplication by $12u^2(e^u-1)^2$ reduces the claim to
\[
(7u^2-12u-12)e^{2u}+2(5u^2+12u+12)e^u-(5u^2+12u+12)\ge0.
\]
Its coefficients through degree four vanish. For $h\ge5$, its coefficient multiplied by $h!$ is
\[
2^{h-2}(7h^2-31h-48)+10h^2+14h+24>0:
\]
at $h=5$ this is $120$, and for $h\ge6$ both displayed terms are positive. This proves \eqref{eq:logmean-curvature}.

Uniformly for the integers $1\le j\le n$,
\begin{equation}\label{eq:eulerian-evaluation}
A_n(j)=n!c(j)^p(1+e_{n,j}),\qquad |e_{n,j}|<2^{-120}.
\end{equation}
The assertion is exact at $j=1$. For $j\ge2$, the Frobenius identity, reciprocity and the partial fractions of $1/(1-e^{-z})$ give
\[
\frac{A_n(j)}{n!c(j)^p}
=\sum_{h\in\mathbb Z}\left(1+\frac{2\pi i h}{\log j}\right)^{-p}.
\]
Put $a=2\pi/\log n$. The nonzero terms have total modulus at most
\begin{equation}\label{eq:evaluation-tail}
2(1+a^2)^{-p/2}
\left(1+\frac1{pa^2/(1+a^2)-1}\right).
\end{equation}
Indeed the decreasing-integral bound applies after using
$1+a^2t^2\ge(1+a^2)t^{2a^2/(1+a^2)}$ for $t\ge1$.
Both $p\log(1+a^2)$ and $pa^2/(1+a^2)$ increase with $n\ge250$.
At $n=250$, $\pi>3$ and $\log250<6$ imply $a^2>1$ and
$pa^2/(1+a^2)>125$. Thus \eqref{eq:evaluation-tail} is less than
$4\cdot2^{-125}<2^{-120}$.

Summing \eqref{eq:bonferroni} against $A(n,h)$ gives
\begin{equation}\label{eq:B-small}
A_n(j)-(n+1)A_n(j-1)\le B(n,j)\le A_n(j).
\end{equation}
For $j\ge2$, \eqref{eq:logmean-curvature} gives
$(c(j-1)/c(j))^p\le e^{-p/(2j)}$; at $j=1$, $B(n,1)=n!$ exactly. Consequently
\begin{equation}\label{eq:B-small-relative}
\left|\frac{B(n,j)}{n!c(j)^p}-1\right|<\frac1{1000}
\qquad(1\le j\le K+4),
\end{equation}
and $0<B(n,j)/(n!c(j)^p)<2$ for every $1\le j\le n$.
For completeness, the error in \eqref{eq:B-small-relative} is at most
\[
2^{-120}+(n+1)(1+2^{-120})
\exp\left(-\frac{n+1}{2(\sqrt n/3+4)}\right).
\]
This decreases for $n\ge250$: after setting $t=\sqrt n$, the relevant logarithmic derivative is
\[
\frac{3(t^2+24t-1)}{2(t+12)^2}-\frac{2t}{t^2+1}>0.
\]
At $n=250$ the exponent exceeds $753/56>13$, and $251e^{-13}<1/1000$ with room for the $2^{-120}$ term.

\subsection{Alternating signs}
For $2\le k\le K+1$ define
\[
r_{n,k}=\left(\frac{c(k-1)}{c(k+1)}\right)^{p/2},
\qquad s_k=-\log r_{n,k}.
\]
Then
\begin{equation}\label{eq:small-separator-sign}
\sgn P_n(-r_{n,k})=(-1)^k.
\end{equation}
To prove this, let $\delta_k=2\ell(k)-\ell(k-1)-\ell(k+1)$.
Equation \eqref{eq:logmean-curvature} yields
$\delta_k\ge5/[12(k+1)^2]$ and $p\delta_k\ge2$.
For $K=5$, use $p\ge251$ and $k+1\le7$; for $K\ge6$, use
$n\ge9K^2$ and $(15/4)(K/(K+2))^2\ge135/64>2$.

Set $r=r_{n,k}$ and $T_j=c(j)^p r^j$. Strict log-concavity gives
\[
\frac{T_{k-h}}{T_k},\frac{T_{k+h}}{T_k}\le q^h,\qquad
q=e^{-p\delta_k/2}\le e^{-1}<\frac38.
\]
Each one-sided alternating tail is decreasing, so
$\sum_{j=1}^n(-1)^{j-k}T_j/T_k\ge1-2q>1/4$.
Replacing the ideal coefficients by $B(n,j)/n!$ changes this by less than
\[
\frac1{1000}\left(1+\frac{2q}{1-q}\right)
+\frac{q^4}{1-q}
<\frac{11}{5000}+\frac{81}{2560}<\frac14.
\]
Here \eqref{eq:B-small-relative} applies through $j=k+3$, and the global ratio in $(0,2)$ applies thereafter. This proves \eqref{eq:small-separator-sign}. The $s_k$ strictly decrease with $k$.

\subsection{Joining these signs to the bulk}
\begin{lemma}\label{lem:small}
There is a phase node $s_*$ with
\[
\phi(s_*)=-(K+1)\pi,\qquad 0<s_*<3\sqrt n,\qquad s_*<s_K.
\]
The phase-node sign followed by the signs at $s_K,s_{K-1},\ldots,s_2$ and near $x=0-$ supplies $K$ distinct negative roots after this node.
\end{lemma}
\begin{proof}
For $s>0$, write $C=1+e^{-s}$, $a_i=C/\beta_i$, $t=C/s$, and
$\mu(t)=tF_n'(t)/F_n(t)$. Directly from \eqref{eq:phase},
\begin{equation}\label{eq:positive-phase-mean}
-\phi(s)=2\arctan(\pi/s)
+\sum_i\arctan\frac{a_i\pi}{s(s+a_i)+\pi^2}
<\frac{\pi}{s}\{\mu(t)+1\}.
\end{equation}
For $0<t\le1$, $\mu(t)\le n-1$. In fact the only negative coefficient of
$(n-1)F_n-tF_n'$ is $-n!t^n$, and the term of degree $n-2$ covers it:
$A(n,3)\ge n(n-1)$ for $n\ge5$. The latter starts at $A(5,3)=66$ and follows inductively from $A(n+1,3)\ge3A(n,3)$. Thus
\[
-\phi(3\sqrt n)<\pi\sqrt n/3<(K+1)\pi.
\]

To compare with $s_K$, give $U>0$ the density proportional to
$g(u)e^{-su}$, where $g(u)=A_n(Cu)$. Termwise integration gives
\[
\E U=\frac{\mu(C/s)+1}{s}=:M.
\]
The function $h(u)=ug'(u)/g(u)$ is concave, since the negative roots of $A_n$ express it as a sum of terms $u/(u+\alpha)$, $\alpha\ge0$. Integration by parts and Jensen's inequality yield
$sM=1+\E h(U)\le1+h(M)$.
As $1/u+g'(u)/g(u)$ strictly decreases, it follows that
\begin{equation}\label{eq:mean-saddle}
s>\frac1b+\frac{g'(b)}{g(b)}\quad\Longrightarrow\quad M<b.
\end{equation}
Take $s=s_K$ and $b=K+1$. The curvature bound gives
\[
s-p\{\ell(b)-\ell(b-1)\}=\frac p2\delta_K\ge1.
\]
Concavity of $\log A_n$ and \eqref{eq:eulerian-evaluation} imply
\[
\frac{g'(b)}{g(b)}
\le C\{\log A_n(b)-\log A_n(b-1)\}
\le C\{p(\ell(b)-\ell(b-1))+4\cdot2^{-120}\}.
\]
Here $b\ge6$, $b\le(2/5)\sqrt n$, $s\ge p/(2b)\ge(5/4)\sqrt n$, and
$p(\ell(b)-\ell(b-1))\le p$. Hence the extra factor $C-1=e^{-s}$ changes the last bound by less than $1/1000$; the function
$(n+1)e^{-(5/4)\sqrt n}$ decreases and is already below $1/1000$ at $250$.
Since $1/b\le1/6$, \eqref{eq:mean-saddle} gives $M<b$. Equation
\eqref{eq:positive-phase-mean} now gives $\phi(s_K)>-(K+1)\pi$.

Finally put $t_0=2/\pi$. Lemma~\ref{lem:weightedmean} and
$\arctan x\ge x/(1+x)$ give
\[
-\phi(0)=\pi+\sum_i\arctan(t_0/\beta_i)
\ge\pi+\mu(t_0)-1
\ge\pi+(n-1)/2>(K+1)\pi.
\]
Strict phase monotonicity proves the three assertions about $s_*$. Its sign is $(-1)^{K+1}$ by Lemma~\ref{lem:bulk}; the subsequent separator signs alternate, and the sign near zero is negative because $B(n,1)=n!$. Their $K$ intervening intervals are disjoint.
\end{proof}
